\documentclass[12pt]{article}
\usepackage{amsmath,amssymb,amsfonts}
\newcommand{\xmapsto}[2][]{\xrightarrow[#1]{#2}}
\title{Theory of Field Extensions Problem set solutions}
\begin{document}
\maketitle
\section{Problem Set 1}
\begin{enumerate}
    \item[\bf{Soln 1:}] $p(x)=x^3 +9x+6$  is irreducible over $\mathbb{Q}$ using Eisenstein's criterion by taking $p=3$ \\
    If $\alpha$ is a root of $p(x)$ to find the inverse of $1+\alpha$ in $\mathbb{Q}(\alpha)$, we note that since $p(x)$ is irreducible, therefore there exist two polynomials $a(x),b(x)  \in F[x]$ such that
\begin{equation}
    a(x)(1+x) + b(x)(x^3+9x+1)=1
\end{equation}


 Since $\alpha$ is a root of $p(x)$ therefore $p(\alpha)=0$ which yields 
 \begin{equation}
    a(\alpha)(1+\alpha)=1 \implies a(\alpha)=(1+\alpha)^{-1}
\end{equation}
So our job is to compute $a(\alpha)$.To do that, we need to use the division algorithm by $1+x$ on $x^3 +9x+6$, where we will get 
\begin{align*}
    &   x^3 +9x+6=(x^2-x+10)(1+x)-4 \\
    &\implies (x^2-x+10)(1+x)-(x^3 +9x+6)=4 \\
    &\implies \frac{1}{4}(x^2-x+10)(1+x)-\frac{1}{4}(x^3 +9x+6)=1
\end{align*}
    which is of the the form $a(x)(1+x) + b(x)(x^3+9x+1)=1$ from (1), where $a(x)=\dfrac{1}{4}(x^2-x+10)$ and $b(x)=-\dfrac{1}{4}$ and thus we have
    
    
    $a(\alpha)=\dfrac{1}{4}(\alpha^2-\alpha +10)=(1+\alpha)^{-1}$ from (2)

    \item[\bf{Soln 2}: ] Let $p(y)=y^3+y+1$. Since $p(y)$ is a cubic polynomial, therefore it suffices to show that there are no roots for $p(y)$ in $\mathbb{F}_2$. We have, $p(0)=1$ and $p(1)=1$ as well. So $p(y)$ does not have any roots in $\mathbb{F}_2$.

    
    For the second part, if $\beta$ is a root of $p(y)$, then 
    \begin{align*}
        &\beta^3+\beta +1=0 \\
     &\implies \beta^3=\beta +1 \tag{1}
    \end{align*} 
    Also, using the relation (1),
    \begin{align*}
        &\beta^4=\beta^3.\beta=(\beta +1)\beta=\beta^2+\beta\\ 
        &\beta^5=\beta^3.\beta^2=(\beta+1)\beta^2=\beta^3+\beta^2=\beta+\beta^2+1\\
        &\beta^6=\beta^3.\beta^3=(\beta +1)(\beta+1)=\beta^2+1 \quad \text{(since $\mathbb{F}_2$ has characteristic 2)} \\
        &\beta^7=\beta.\beta^6=\beta(\beta^2+1)=\beta^3+\beta =\beta+1+\beta=2\beta+1=1 \quad \text{(Since $2\beta=0$ in $\mathbb{F}_2(\beta)$ )}
    \end{align*}

    Therefore, in $\mathbb{F}_2(\beta)$, the different powers of $\beta$ are $1, \beta, \beta^2,\beta+1,\beta^2+\beta$ and $\beta^2+\beta+1$
    \item[\bf{Soln 3:} ] Suppose $f(x)=x^3-nx+2 $ is reducible over $\mathbb{Q}$ and let $\alpha \in \mathbb{Q}$ be a root of $f(x)$,i.e $f(\alpha)=0\\
    \implies \alpha^3-n\alpha+2=0 \\
    \implies \alpha(\alpha^2-n)=-2 \\
    \implies \alpha|-2
    \implies \alpha=1,-1,2-2$

    For $\alpha=-1$, we have
    $f(-1)= (-1)+n+2=1+n$
    which means that for $f(-1)=0$, we have $n=-1$.
    Similarly, $f(1)=0
    \implies 3-n=0
    \implies n=3$


    $f(2)=0
    \implies 10-2n=0
    \implies n=5$


    $f(-2)=0
    \implies -6+2n=0
    \implies n=3$

    
    Thus for $f(x)=0$, that is for $f(x)$ to be reducible over $\mathbb{Q}$, we must have $n=-1,3$ or $5$

    \item[\bf{Soln 4}: ] The inclusion $\mathbb{Q}(2 + \sqrt3) \subseteq \mathbb{Q}(\sqrt3)$ is obvious. For the reverse inclusion, $\sqrt3=2-2+\sqrt3 =(2+\sqrt3)-2\in \mathbb{Q}(2+\sqrt3)$. Thus $\mathbb{Q}(\sqrt3)=\mathbb{Q}(2+\sqrt3)$ and therefore $[\mathbb{Q}(2+\sqrt3):\mathbb{Q]}=2$
    \item[\bf{Soln 5}: ] Let $\alpha=1+i$
    $\begin{aligned}[t]
        &\implies \alpha-1=i\\
        &\implies \alpha^2-2\alpha+1=-1\\ 
        &\implies \alpha^2-2\alpha+2=0
    \end{aligned}$\\
    Therefore $p(x)=x^2 - 2x+2$ is the required minimal polynomial.
    \item[\bf{Soln 6}: ] True since if $\alpha$ is a root of $y^3-2$, then $\alpha^3 -2=0 \implies \alpha = 2^{\frac{1}{3}} \notin \mathbb{Q}(i)$. Same thing applies to $y^3-3$
    
    \item[\bf{Soln 7}: ] Let $p(x)$ be an irreducible factor of $h(g(x))$.\\
    Let deg($p(x))=n$. We have to show $m|n$.

    Let $\alpha$ be a root of $p(x)$. Now 
    \[\begin{aligned}
        &p(x)|h(g(x))\\
    &\implies h(g(x))=p(x)b(x) \quad \text{for some $b(x) \in F[x]$}\\
    &\implies h(g(\alpha))=p(\alpha)b(\alpha)=0 \\
    &\implies g(\alpha) \quad \text{is a root of $h(g(x))$}    
    \end{aligned}\]
     Thus, we have the tower of fields 
     \[F \subset F(g(\alpha))\subset F(\alpha)\]\\
     Here, $g(\alpha) \in F(\alpha)$ and therefore $F(g(\alpha)) \subset F(\alpha)$.\\
    Using the tower law, we have
\[
\begin{aligned}
n 
&=   [F(\alpha):F] \\
&=   [F(\alpha):F(g(\alpha))]\,[F(g(\alpha)):F] \\
&=         [F(\alpha):F(g(\alpha))]\, m \\
\implies& n = [F(\alpha):F(g(\alpha))]\, m \\
\implies& m \mid n.
\end{aligned}
\]
    \item[\bf{Soln 8}: ]  Denote this map by $\varphi$. Then
    \[
    \varphi(a + b\sqrt{2} + c + d\sqrt{2})= \varphi((a +c)+ (b+d)\sqrt2)= a + c - b\sqrt{2} - d\sqrt{2} = \varphi(a + b\sqrt{2}) + \varphi(c + d\sqrt{2})\] and
    \[
    \begin{aligned}
    \varphi\!\left((a + b\sqrt{2})(c + d\sqrt{2})\right)
    &= \varphi(ac + 2bd + (ad + bc)\sqrt{2}) \\
    &= ac + 2bd - (ad + bc)\sqrt{2} \\
    &=ac+2bd-ad\sqrt2-bc\sqrt2\\
    &=ac- ad \sqrt2+2bd-bc\sqrt2 \\
    &=a(c-d\sqrt2)-b\sqrt2(c-d\sqrt2)\\
    &= (a - b\sqrt{2})(c - d\sqrt{2}) \\
    &= \varphi(a + b\sqrt{2})\,\varphi(c + d\sqrt{2}),
    \end{aligned}
    \]
    so $\varphi$ is a homomorphism. If $\varphi(a + b\sqrt{2}) = \varphi(c + d\sqrt{2})$, then
    \[
    a - b\sqrt{2} = c - d\sqrt{2}.
    \]
    \[
    \implies a-c=(b-d)\sqrt2
    \]
     If $b-d \neq 0$, then $\dfrac{a-c}{b-d}=\sqrt2$, which is a contradiction, so \\ $b-d=0\implies a-c=0$ and hence $a = c$ and $b = d$. Thus $\varphi$ is injective.  
    Furthermore, given $a + b\sqrt{2} \in \mathbb{Q}(\sqrt{2})$, we have
    \[
    \varphi(a - b\sqrt{2}) = a + b\sqrt{2},
        \]
    so $\varphi$ is surjective. Therefore $\varphi$ is an automorphism of $\mathbb{Q}(\sqrt{2})$.
    \item[\bf{Soln 9}: ] Since $\sqrt{2} + \sqrt{3}$ is an element of $\mathbb{Q}(\sqrt{2}, \sqrt{3})$, clearly
    \[\mathbb{Q}(\sqrt{2} + \sqrt{3}) \subset \mathbb{Q}(\sqrt{2}, \sqrt{3}).\]
    On the other hand, consider $\theta = \sqrt{2} + \sqrt{3}$. Then
    \[
    \theta^2 = 5 + 2\sqrt{6}, \qquad \theta^3 = 11\sqrt{2} + 9\sqrt{3},
    \]
    so
    \[
    \sqrt{2} = \frac{1}{2}(\theta^3 - 9\theta), \qquad 
    \sqrt{3} = \frac{1}{2}(11\theta - \theta^3).
    \]
    Therefore $\sqrt{2} \in \mathbb{Q}(\theta)$ and $\sqrt{3} \in \mathbb{Q}(\theta)$, so
    \[
    \mathbb{Q}(\sqrt{2},\sqrt{3}) \subset \mathbb{Q}(\sqrt{2} + \sqrt{3}).
    \]
    It follows that
    \[
    \mathbb{Q}(\sqrt{2} + \sqrt{3}) = \mathbb{Q}(\sqrt{2}, \sqrt{3}),
    \]
    so that
    \[
    [\mathbb{Q}(\sqrt{2} + \sqrt{3}) : \mathbb{Q}] 
    = [\mathbb{Q}(\sqrt{2}, \sqrt{3}) : \mathbb{Q}]
    = 4.
    \]

    We also have
    \[
    \theta^4 - 10\theta^2 
    = (49 + 20\sqrt{6}) - 10(5 + 2\sqrt{6})
    = -1,
    \]
    so
    \[
    \theta^4 - 10\theta^2 + 1 = 0.
    \]

    Since $[\mathbb{Q}(\sqrt{2} + \sqrt{3}) : \mathbb{Q}] = 4$, the polynomial
    \[
    x^4 - 10x^2 + 1
    \]
    is irreducible over $\mathbb{Q}$, and is satisfied by $\sqrt{2} + \sqrt{3}$

    \item[\bf{Soln 11}: ] Given, $\beta_n^2 \in \mathbb{Q}$\\
     $\implies \beta_n^2 \in \mathbb{Q}(\beta_1,\beta_2,...,\beta_{n-1})$. But $\beta_n \notin \mathbb{Q}(\beta_1,\beta_2,...,\beta_{n-1})$. Therefore ,$x^2-\beta_n^2$ is the minimal polynomial of $\beta_n$ over $\mathbb{Q}(\beta_1,\beta_2,...,\beta_{n-1})$.\\
     This means that \[[\mathbb{Q}(\beta_1,\beta_2,...,\beta_n):\mathbb{Q}(\beta_1,\beta_2,...,\beta_{n-1})]=2.\]
     

     Therefore,\[
     \begin{aligned}
     &[\mathbb{Q}(\beta_1,\beta_2,...,\beta_n): \mathbb{Q}] \\ 
     & =[\mathbb{Q}(\beta_1,\beta_2,...,\beta_n):\mathbb{Q}(\beta_1,\beta_2,...,\beta_{n-1})] \\ &\times  [\mathbb{Q}(\beta_1,\beta_2,...,\beta_{n-1}):\mathbb{Q}(\beta_1,\beta_2,...,\beta_{n-2})] \times ...  \times [\mathbb{Q}(\beta_1): \mathbb{Q}]=2^n\\ 
      \end{aligned}
      \]

     Now suppose $2^\frac{1}{3} \in F$, i.e $[\mathbb{Q}(2^\frac{1}{3}):\mathbb{Q}]=3$\\
     We have 

     $\mathbb{Q} \subset \mathbb{Q}(2^\frac{1}{3}) \subset F$, so again using the Tower Law,
     \[
     \begin{aligned}
     &[F:\mathbb{Q}]= [F:\mathbb{Q}(2^\frac{1}{3})]\times[\mathbb{Q}(2^\frac{1}{3}:\mathbb{Q}] \\
     &\implies 2^n=[F:\mathbb{Q}(2^\frac{1}{3})] \times 3 \\
     &\implies \dfrac{2^n}{3}=[F:\mathbb{Q}(2^\frac{1}{3})]
      \end{aligned}
      \]
      which is a contradiction.
      So, $2^\frac{1}{3} \notin F$ 
      \item[\bf{Soln 12: }] Suppose $\sqrt{a}+\sqrt{b}=\sqrt{m}+\sqrt{n}$ for some $m,n\in F$, so that 
\[
a+\sqrt{b}=m+n+2\sqrt{mn}.
\]
Since $b$ is not a square in $F$, this means $\sqrt{b}=2\sqrt{mn}$. We also have 
\[
\sqrt{a}+\sqrt{b}-\sqrt{n}=\sqrt{m},
\]
so
\[
\sqrt{b}=2\sqrt{n}\bigl(\sqrt{a}+\sqrt{b}-\sqrt{n}\bigr).
\]
Hence,
\[
\sqrt{b}=2\sqrt{n}(a+\sqrt{b})-2n
\]
\[
\Rightarrow\qquad (\sqrt{b}+2n)^2 = 4n(a+\sqrt{b})
\]
\[
\Rightarrow\qquad b + 4n\sqrt{b} + 4n^2 = 4n(a+\sqrt{b})
\]
\[
\Rightarrow\qquad b + 4n^2 - 4na = 0
\]
\[
\Rightarrow\qquad n = \frac{4a \pm \sqrt{16a^2 - 16b}}{8}
= \frac{a \pm \sqrt{a^2 - b}}{2}.
\]
Therefore, since $a$ and $n$ belong to $F$, so does $\sqrt{a^2 - b}$.

Conversely, assume that $a^2-b$ is a square in $F$, so that $\sqrt{a^2-b}\in F$. We prove that there exist
$m,n\in F$ such that 
\[
\sqrt{a}+\sqrt{b}=\sqrt{m}+\sqrt{n}.
\]
Consider
\[
m=\frac{a+\sqrt{a^2-b}}{2}
\qquad\text{and}\qquad
n=\frac{a-\sqrt{a^2-b}}{2}.
\]

Note that $m$ and $n$ belong to $F$ as $\operatorname{char}(F)\neq 2$. We claim 
\[
\sqrt{a}+\sqrt{b} = \sqrt{m}+\sqrt{n}.
\]
Indeed, we have
\[
m = \frac{(a+\sqrt{b}) + 2\sqrt{a^2-b} + (a-\sqrt{b})}{4}
= \left(\frac{\sqrt{a}+\sqrt{b} + \sqrt{a}-\sqrt{b}}{2}\right)^2.
\]
and
\[
n=\frac{(a+\sqrt{b}) - 2\sqrt{a^{2}-b} + (a-\sqrt{b})}{4}
= \left( \frac{\sqrt{a}+\sqrt{b} - \sqrt{a}-\sqrt{b}}{2} \right)^{2}.
\]

Thus
\[
\sqrt{m}=\frac{\sqrt{a}+\sqrt{b}+\sqrt{a}-\sqrt{b}}{2}, 
\qquad 
\sqrt{n}=\frac{\sqrt{a}+\sqrt{b}-\sqrt{a}+\sqrt{b}}{2},
\]
so
\[
\sqrt{m}+\sqrt{n}
= \frac{\sqrt{a}+\sqrt{b}+\sqrt{a}-\sqrt{b}}{2}
+ \frac{\sqrt{a}+\sqrt{b}-\sqrt{a}+\sqrt{b}}{2}
= \sqrt{a}+\sqrt{b},
\]
as claimed.

Finally, we use this to determine when the field $\mathbb{Q}(\sqrt{a}+\sqrt{b})$,  
$a,b\in\mathbb{Q}$, is biquadratic over $\mathbb{Q}$.

If $a^2-b$ is a square in $\mathbb{Q}$ and $b$ is not, we have
\[
\mathbb{Q}(\sqrt{a}+\sqrt{b}) = \mathbb{Q}(\sqrt{m}+\sqrt{n})
= \mathbb{Q}(\sqrt{m},\sqrt{n}),
\]
so, $\mathbb{Q}(\sqrt{a}+\sqrt{b})$ is biquadratic over $\mathbb{Q}$
when $a^{2}-b$ is a square in $\mathbb{Q}$ and none of $b$, $m$, $n$ or $mn$
is a square in $\mathbb{Q}$.

Since
\[
mn = \frac{a+\sqrt{a^{2}-b}}{2}\cdot \frac{a-\sqrt{a^{2}-b}}{2}
= \frac{a^{2}-(a^{2}-b)}{4}
= \frac{b}{4},
\]
$mn$ is never a square when $b$ is not. Thus $\mathbb{Q}(\sqrt{a}+\sqrt{b})$
is biquadratic over $\mathbb{Q}$ exactly when $a^{2}-b$ is a square in $\mathbb{Q}$
and none of $b$, $m$, or $n$ is a square in $\mathbb{Q}$.


    \end{enumerate}

\section{Problem Set 2}
\begin{enumerate}
    \item[\bf{Soln 1: }] $[\mathbb{Q}(\pi):\mathbb{Q}(\pi^3)]=3$ since deg($x^3-\pi^3$)=3 
    \item[\bf{Soln 2: }]  $[\mathbb{Q}(\sqrt5,\sqrt7):\mathbb{Q}]=4$ 

    
    Justification:  On adjoining $\sqrt5$ to $\mathbb{Q}$ we get a quadratic extension, i.e., $[\mathbb{Q}(\sqrt5):\mathbb{Q}]=2$ since deg$(x^2-5)=2$. Since $\sqrt7 \notin \mathbb{Q}(\sqrt5)$, therefore on adjoining $\sqrt7$ to $\mathbb{Q}(\sqrt5)$, we get another extension of degree 2, that is, $[\mathbb{Q}(\sqrt5,\sqrt7):\mathbb{Q}(\sqrt5)]=2$. Hence on applying the tower law, we get
    \[
        \begin{aligned}
            [\mathbb{Q}(\sqrt5,\sqrt7):\mathbb{Q}]=&[\mathbb{Q}(\sqrt5,\sqrt7):\mathbb{Q}(\sqrt5)] \times [\mathbb{Q}(\sqrt5):\mathbb{Q}]\\
            =&2 \times 2=4
        \end{aligned}
    \]
    \item[\bf{Soln 3: }] $[\mathbb{Q}(\sqrt2,\sqrt5,\sqrt7:\mathbb{Q})]=8$ using similar reasoning as in solution 2
    \item[\bf{Soln 7 :} ]  True. The polynomials $x^2-2$ and $x^2 -8$ have the same splitting field $\mathbb{Q}(\sqrt2)$. One more example: the polynomials $x^4+2$ and $x^4-2$ which have the same splitting field $\mathbb{Q}(2^\frac{1}{4},i)$
    \item[\bf{Soln 8}: ] Same as in soln 7 no?
    \item[\bf{Soln 9a,b and c}] False. Take $\mathbb{Q}(2^{1/3})$ over $\mathbb{Q}$. 
    \item[\bf{Soln 10: }] Consider the field $E=\mathbb{F}_p(t)$ where $t$ is an indeterminate. The Frobenius map $x \longmapsto x^p$ is not surjective.
    \item[\bf{Soln 11a: }] False. Let us take the map $\theta : \mathbb{C \longrightarrow C}$ defined by $\theta(a+ib)=a-ib.$ Here we are considering $\mathbb{C}$ as an extension of $F=\mathbb{Q}(i)$. So the map $\theta$ is a field isomorphism but not an $F$-isomorphism, since if $x_1,y_1\in \mathbb{Q}$, then $\theta(x_1+iy_1)=x_1-iy_1$. That is, $\theta$ does not fix the elements of $\mathbb{Q}(i)$
    \item[\bf{Soln 12: }] False. Suppose $\sqrt2$ and $\sqrt3$ are conjugates and consider the map $\phi: \mathbb{Q}(\sqrt2) \longrightarrow\mathbb{Q}(\sqrt2)$ where $\phi(\sqrt2)=\sqrt3$, since $\sqrt2$ and $\sqrt3$ are conjugates. Now,
    \[
    \begin{aligned}
        &\phi(2)=\phi(\sqrt2.\sqrt2)\\
        &\implies2=\phi(\sqrt2)\phi(\sqrt2) \qquad \text{(Since $\phi$ fixes the elements of $\mathbb{Q}$ and is also a homomorphism)}\\
        &\implies2=\sqrt3\times\sqrt3\\
        &\implies2=3
    \end{aligned}
    \]
    which is a contradiction, so $\sqrt2$ and $\sqrt3$ are not conjugates over $\mathbb{Q}$
    \item[\bf{Soln 13: }] We have $p(x)=x^3+x^2+1$ over $\mathbb{Z}_2$. Here $p(x)$ is irreducible over $\mathbb{Z}_2$ since it is a cubic and it does not have any roots in $\mathbb{Z}_2$. Now, if $\alpha$ is a root of $p(x)$, that is, is $p(\alpha)=0\implies \alpha^3+\alpha^2+1=0\implies\alpha ^3=\alpha^2+1$


    Now we want to check if $\alpha^2$ is a root. So, we evaluate $p(\alpha^2).$

    We have,
    \begin{equation*}
        p(\alpha^2)=\alpha^6+\alpha^4+1 \tag{1}
    \end{equation*}

    Here, we compute the $\alpha^6$ and $\alpha^4$. So, $\alpha^4=\alpha^3.\alpha=(\alpha^2+1)\alpha=\alpha^3+\alpha=\alpha^2+\alpha+1$ and $\alpha^6=(\alpha^3)^2=(\alpha^2+1)^2=\alpha^4+1=\alpha^3+\alpha+1$, since $\mathbb{Z}_2$ has characteristic 2. So substituting in (1), we have,
    \[
    \begin{aligned}
        p(\alpha^2)=&\alpha^3+\alpha+1+\alpha^3+\alpha+1\\
        =&2\alpha^3+2\alpha+2=0 \qquad \text{(Again since char $\mathbb{Z}_2=2$)}
\end{aligned}
    \]

    Now, if $\beta$ is a third root, then $p(x)=(x-\alpha)(x-\alpha^2)(x-\beta)$. That is, $x^3+x^2+1=(x-\alpha)(x-\alpha^2)(x-\beta)$. Comparing the coefficients of $x^2$, we get that 
    \[
    1=-\alpha-\alpha^2-\beta\implies\beta=1+\alpha+\alpha^2
    \] 

    Therefore the splitting field of $p(x)$ is $\mathbb{Z}_2(\alpha,\alpha^2,1+\alpha+\alpha^2)=\mathbb{Z}_2(\alpha)$

    Therefore, we have
    \[
    \mathbb{Z}_2(\alpha) \cong \mathbb{Z}_2[x]/<p(x)>
    \]
\end{enumerate}

\section{Problem Set 3 }

    \begin{enumerate}
        \item[\bf{Soln 1: }]  The degree of the splitting field of $x^4-2$ over $\mathbb{Q}$ is 8. 


        Reason: The splitting field of $x^4-2$ over $\mathbb{Q}$ is $\mathbb{Q}(2^\frac{1}{4},i)$. Therefore the degree is 
         \[   
        \begin{aligned}
            [\mathbb{Q}(2^\frac{1}{4},i):\mathbb{Q}]&= [\mathbb{Q}(2^\frac{1}{4},i):\mathbb{Q}(i)]\times[\mathbb{Q}(i):\mathbb{Q}]\\
            &=4\times2=8
        \end{aligned}
        \]

        \item[\bf{Soln 2: }] The degree of the splitting field of $x^3-2$ over $\mathbb{Q}$ is 6.

        Reason: The splitting field of $x^3-2$ over $\mathbb{Q}$ is $\mathbb{Q}(2^\frac{1}{3},\omega)$, where $\omega$ is a cube root of unity. 

        Therefore the degree is 
        \[
        \begin{aligned}
            [\mathbb{Q}(2^\frac{1}{3},\omega):\mathbb{Q}]&=[\mathbb{Q}(2^\frac{1}{3},\omega):\mathbb{Q}(\omega)] \times [\mathbb{Q}(\omega):\mathbb{Q}]\\
            &=3\times2=6
        \end{aligned}
        \]

        \item[\bf{Soln 3: }] We have \[
        \begin{aligned}
        x^4+x^2+1=(x^2+1)^2-x^2&=(x^2+x+1)(x^2-x+1)=0 \\
                                \implies x=\dfrac{-1}{2}\pm\dfrac{\sqrt3}{2}i,\dfrac{1}{2}\pm\dfrac{\sqrt3}{2}i
    \end{aligned}   
        \]
        Therefore the splitting field of $x^4+x^2+1$ is $\mathbb{Q}(\sqrt{-3})$ or $\mathbb{Q}(\sqrt{3}i)$
    \item[\bf{Soln 4: }] 
    \item[\bf{Soln 5: }] The roots of $x^4+4$ are 

    \[ x^4+4=0\implies (x^2+2x+2)(x^2-2x+2)=0\implies x=1\pm i,-1\pm i\].

    So the splitting field of $x^4+4$ over $\mathbb{Q}$ is $\mathbb{Q}(i)$

    \item[\bf{}] 
    \end{enumerate}

\section{Problem Set 4}
\begin{enumerate}
    \item[\bf{Soln 1 and 2: }] To recall the definition of an algebraic closure, it is a field $\bar F$ such that
    \begin{enumerate}
        \item $F \subseteq \bar F$ and $\bar F/F $ is algebraic.
        \item $\bar F$ is algebraically closed.
        \item $\bar F$ is the smallest field satsifying both conditions (a) and (b)
    \end{enumerate}

    For condition (a), we have  $\overline{\mathbb{Q}}={\mathbb{A}}=\{\alpha \in \mathbb{C} \mid  \text{$\alpha$ is algebraic over $\mathbb{Q}$}\}$. So $\mathbb{Q} \subset \mathbb{A}$ is obvious and $\mathbb{A/\mathbb{Q}}$ is algebraic by definition. So condition (a) is satisfied.


    For condition (b), let us take a non-constant polynomial $f(x)$ over $\mathbb{A}$. Then $f(x)$ has its roots in some splitting field $E$ of $f(x)$. Suppose $\alpha$ is a root of $f(x)$. Then $\alpha$ is algebraic over $\mathbb{A}$. But $\mathbb{A}$ is algebraic over $\mathbb{Q}$, therefore we obtain that $\alpha$ is algebraic over $\mathbb{Q}$ as well. Hence, $\alpha \in \mathbb{A}$ proving that all the roots of $f(x)$ are in $\mathbb{A}$.

    Now to prove condition (c), let $E$ be an algebraic extension of $\mathbb{Q}$ satisfying conditions (a) and (b). If $\beta \in \mathbb{A}$, then by definition, $\beta$ is algebraic over $\mathbb{Q}$. It then follows that $\beta \in E$, since $E$ is algebraically closed. So $\mathbb{A} \subset E$, so $\mathbb{A}$ is the smallest field which satisfies (a) and (b). Thus condition (c) is proved

    \item[\bf{Soln 3: }] False. Consider the polynomial $p(y)=y^2-x$ over $\mathbb{C}(x)$. That is $p(y) \in \mathbb{C}(x)[y]$, where $y$ is an indeterminate. The roots of $p(y)$ are $\pm \sqrt{x}$, which does not exist in $\mathbb{C}(x)$. 
    \item[\bf{Soln 4: }] False. We know that the algebraic closure always exists for any field always exists. So take the algebraic closure of $\mathbb{Z}_p$,i.e take $\overline{\mathbb{Z}}_p$. This field is of characteristic $p$.
    \item[\bf{Soln 6: }] Take the extension $\mathbb{Q}(\sqrt2)$ over $\mathbb{Q}$. $\mathbb{Q}(\sqrt2)$ is normal over $\mathbb{Q}$ since it is the splitting field of $x^2-2$ over $\mathbb{Q}$, but it is not the algebraic closure of $\mathbb{Q}$
    \item[\bf{Soln 7: }] False. Take the extension $\mathbb{C}/\mathbb{Q}$. This extension is normal but not finite.
    \item[\bf{Soln 8: }] False. Take the extension $\mathbb{Q}(2^\frac{1}{3})/\mathbb{Q}$. This extension is finite but it is not normal since it does not contain all the roots of the polynomial $x^3-2$ over $\mathbb{Q}$
    \item[\bf{Soln 9: }] False. Take $\mathbb{C}/\mathbb{Q}$, which is normal but not algebraic since there exists an element in $\mathbb{C}$ which is not algebraic over $\mathbb{Q}$, namely $\pi$.
    \item[\bf{Soln 10: }] False. Take the same example in solution 8.
    \item[\bf{Soln 12: }] Take the same extensions as in solutions 7 and 9
\end{enumerate}

\section{Problem Set 5}
\begin{enumerate}
    \item[\bf{Soln 1: }] Let $\alpha=\sqrt{2+\sqrt2}$\\Then
    $\alpha^2=2+\sqrt2\implies \alpha^2-2=\sqrt2\implies\alpha^4-2\alpha+2=0$
    Therefore the minimal polynomial of $\sqrt{2+\sqrt2}$ is $p(x)=x^4-4x^2+2$, which is irreducible over $\mathbb{Q}$ by Eisenstein, and whose roots are distinct(check) and therefore $p(x)$ is separable.
    Also the splitting field of $\sqrt{2+\sqrt2}$ is $\mathbb{Q}(\sqrt{2+\sqrt2})$, which shows that $\mathbb{Q}(\sqrt{2+\sqrt2})/\mathbb{Q}$ is a Galois extension. It is also an extension of degree 4 since $p(x)$ is of degree 4.
    \item[\bf{Soln 2: }] Given the polynomial $p(x)=x^4-14x^2+9$. 
    Putting $t=x^2$, the polynomial transforms into a quadratic $t^2-14t+9$ where we get the roots to be\\
    \[
    \begin{aligned}
    &t=7\pm2\sqrt{10}\\
    &\implies x^2=7\pm2\sqrt{10}\\
    &\implies x=\sqrt{7\pm2\sqrt{10}}
    \end{aligned}
    \]

    Here we see that $7\pm2\sqrt{10}=(\sqrt5\pm\sqrt2)^2$.
     So the roots can be written as $x=\pm(\sqrt5\pm\sqrt2)$. Thus the splitting field of this polynomial is $\mathbb{Q}(\sqrt2+\sqrt5)=\mathbb{Q}(\sqrt2,\sqrt5)$

     The Galois group of a polynomial is defined to be all the automorphisms of the splitting field which fix the elements of the base field. That is
     
     
     $Gal(\mathbb{Q}(\sqrt2,\sqrt5)/\mathbb{Q})=\{\sigma \in Aut_\mathbb{Q}(\mathbb{Q}(\sqrt2,\sqrt5))\mid \sigma\vert_\mathbb{Q}=1_\mathbb{Q}\}$. So we have to determine each of these maps.

        \begin{enumerate}
            \item[1.] $\sigma_1=Id$
            \item[2.] $\sigma_2:
\begin{array}{rcl}
\sqrt{2} &\longmapsto& -\sqrt{2},\\[2pt]
\sqrt{5} &\longmapsto& \sqrt{5}.
\end{array}$
            \item[3.] $\sigma_3:
\begin{array}{rcl}
\sqrt{2} &\longmapsto& -\sqrt{2},\\[2pt]
\sqrt{5} &\longmapsto& -\sqrt{5}.
\end{array}$
        \item[4.] $\sigma_4:
\begin{array}{rcl}
\sqrt{2} &\longmapsto& \sqrt{2},\\[2pt]
\sqrt{5} &\longmapsto& -\sqrt{5}.
\end{array}$
        \end{enumerate}

    So we have the set $\{Id=\sigma_1, \sigma_2,\sigma_3,\sigma_4\}$ to be the Galois group of $\mathbb{Q}(\sqrt2,\sqrt5)$ over $\mathbb{Q}$.
    \item[\bf{Soln 3: }]  We see that $\mathbb{Q}(\sqrt2,\sqrt5)$ is the splitting field of a separable polynomial $(x^2-2)(x^2-5)$ over $\mathbb{Q}$. So $\mathbb{Q}(\sqrt2,\sqrt5)$ is a Galois extension. In Soln 2, we have found the Galois group of $\mathbb{Q}(\sqrt2,\sqrt5)$ to be
    $\{Id=\sigma_1, \sigma_2,\sigma_3,\sigma_4\}$. To show that it is isomorphic to the Klein 4-group, we note that $\sigma_2^2(\sqrt2)=\sigma_2(\sigma_2(\sqrt2)=\sigma_2(-\sqrt2)=-\sigma_2(\sqrt2)=-(-\sqrt2)=\sqrt2$. That is $\sigma_2^2=Id$. So the order of $\sigma_2$ is 2. We can perform a similar calculation to see that the order of $\sigma_3$ and $\sigma_4$ is also 2. That is the order of all the non-identity elements is 2. So $Gal(\mathbb{Q}(\sqrt2,\sqrt5)/\mathbb{Q}) \cong V_4$, the Klein's 4-group

    \item[\bf{Soln 4: }]  The roots of the polynomial $x^8-3$ are $3^{1/8},3^{1/8}\omega,3^{1/8} \omega^2,...,3^{1/8}\omega^7$, where $\omega$ is the 8-th root of unity. That is $\omega=e^\frac{2\pi i}{8}=e^\frac{\pi i}{4}=\dfrac{1+i}{\sqrt2}.$ So we see that $\omega^2=i,\omega^3=\dfrac{-1+i}{\sqrt2},\omega^4=-1,\omega^5=\dfrac{-1-i}{\sqrt2}, \omega^6=-i,\omega^7\dfrac{1-i}{\sqrt2}=$. So the roots of $x^8-3$ can also be written as $3^{1/8},3^{1/8}\bigg(\dfrac{1+i}{\sqrt2}\bigg),3^{1/8}i,3^{1/8}\bigg(\dfrac{-1+i}{\sqrt2}\bigg),\\-3^{1/8},3^{1/8}\bigg(\dfrac{-1-i}{\sqrt2}\bigg),-3^{1/8}i,3^{1/8}\bigg(\dfrac{1-i}{\sqrt2}\bigg)$. So the splitting field of $x^8-3$ is $\mathbb{Q}(3^{1/8},\omega)$ or $\mathbb{Q}(3^{1/8},\sqrt2,i)$


    Now, to find the permutations of the roots, we must note that every automorphism $\omega \longmapsto\omega^r$ is possible only if $r\in \{1,3,5,7\}$. For if $r$ is one of the even integers modulo 8, say 2, then for example the map $\omega \longmapsto\omega^2=i$ sends $\omega$ to $i$, which is not a primitive(\textit{recall defn of primitive root of unity}) 8th root of unity. So, if $\sigma$ is a permutation of $\omega$, then it must send $\omega$ to a primitive 8th root of unity only. That is $\sigma(\omega)$ must be a primitve root of unity, which is possible only if $r=1,3,5$ or 7. As for $3^{1/8}$, it can go to any other root of the same polynomial, which includes $3^{1/8}\omega^2$, etc. With that being said , we can now determine the individual permutations of the roots: 
    \begin{enumerate}
        \item[1. ] $\sigma_1=id$
        \item[2. ] $\sigma_2: \omega \xmapsto{} \omega,\qquad 3^{1/8} \xmapsto{} 3^{1/8}\omega$
        \item[3. ] $\sigma_3: \omega \xmapsto{} \omega,\qquad 3^{1/8} \xmapsto{} 3^{1/8}\omega^2$
        \item[4. ] $\sigma_4: \omega \xmapsto{} \omega,\qquad 3^{1/8} \xmapsto{} 3^{1/8}\omega^3$
        \item[5. ] $\sigma_5: \omega \xmapsto{} \omega,\qquad 3^{1/8} \xmapsto{} 3^{1/8}\omega^4$
        \item[6. ] $\sigma_6: \omega \xmapsto{} \omega,\qquad 3^{1/8} \xmapsto{} 3^{1/8}\omega^5$
        \item[7. ] $\sigma_7: \omega \xmapsto{} \omega,\qquad 3^{1/8} \xmapsto{} 3^{1/8}\omega^6$
        \item[8. ] $\sigma_8: \omega \xmapsto{} \omega,\qquad 3^{1/8} \xmapsto{} 3^{1/8}\omega^7$
        \item[9. ] $\sigma_9: \omega \xmapsto{} \omega^3,\qquad 3^{1/8} \xmapsto{} 3^{1/8}$
        \item[10. ] $\sigma_{10}: \omega \xmapsto{} \omega^3,\qquad 3^{1/8} \xmapsto{} 3^{1/8}\omega$
        \item[11. ] $\sigma_{11}: \omega \xmapsto{} \omega^3,\qquad 3^{1/8} \xmapsto{} 3^{1/8}\omega^2$
        \item[12. ] $\sigma_{12}: \omega \xmapsto{} \omega^3,\qquad 3^{1/8} \xmapsto{} 3^{1/8}\omega^3$
        \item[13. ] $\sigma_{13}: \omega \xmapsto{} \omega^3,\qquad 3^{1/8} \xmapsto{} 3^{1/8}\omega^4$
        \item[14. ] $\sigma_{14}: \omega \xmapsto{} \omega^3,\qquad 3^{1/8} \xmapsto{} 3^{1/8}\omega^5$
        \item[15. ] $\sigma_{15}: \omega \xmapsto{} \omega^3,\qquad 3^{1/8} \xmapsto{} 3^{1/8}\omega^6$
        \item[16. ] $\sigma_{16}: \omega \xmapsto{} \omega^3,\qquad 3^{1/8} \xmapsto{} 3^{1/8}\omega^7$
        \item[17. ] $\sigma_{17}: \omega \xmapsto{} \omega^5,\qquad 3^{1/8} \xmapsto{} 3^{1/8}$
        \item[18. ] $\sigma_{18}: \omega \xmapsto{} \omega^5,\qquad 3^{1/8} \xmapsto{} 3^{1/8}\omega$
        \item[19. ] $\sigma_{19}: \omega \xmapsto{} \omega^5,\qquad 3^{1/8} \xmapsto{} 3^{1/8}\omega^2$
        \item[20. ] $\sigma_{20}: \omega \xmapsto{} \omega^5,\qquad 3^{1/8} \xmapsto{} 3^{1/8}\omega^3$
        \item[21. ] $\sigma_{21}: \omega \xmapsto{} \omega^5,\qquad 3^{1/8} \xmapsto{} 3^{1/8}\omega^4$
        \item[22. ] $\sigma_{22}: \omega \xmapsto{} \omega^5,\qquad 3^{1/8} \xmapsto{} 3^{1/8}\omega^5$
        \item[23. ] $\sigma_{23}: \omega \xmapsto{} \omega^5,\qquad 3^{1/8} \xmapsto{} 3^{1/8}\omega^6$
        \item[24. ] $\sigma_{24}: \omega \xmapsto{} \omega^5,\qquad 3^{1/8} \xmapsto{} 3^{1/8}\omega^7$
        \item[25. ] $\sigma_{25}: \omega \xmapsto{} \omega^7,\qquad 3^{1/8} \xmapsto{} 3^{1/8}$
        \item[26. ] $\sigma_{26}: \omega \xmapsto{} \omega^7,\qquad 3^{1/8} \xmapsto{} 3^{1/8}\omega$
        \item[27. ] $\sigma_{27}: \omega \xmapsto{} \omega^7,\qquad 3^{1/8} \xmapsto{} 3^{1/8}\omega^2$
        \item[28. ] $\sigma_{28}: \omega \xmapsto{} \omega^7,\qquad 3^{1/8} \xmapsto{} 3^{1/8}\omega^3$
        \item[29. ] $\sigma_{29}: \omega \xmapsto{} \omega^7,\qquad 3^{1/8} \xmapsto{} 3^{1/8}\omega^4$
        \item[30. ] $\sigma_{30}: \omega \xmapsto{} \omega^7,\qquad 3^{1/8} \xmapsto{} 3^{1/8}\omega^5$
        \item[31. ] $\sigma_{31}: \omega \xmapsto{} \omega^7,\qquad 3^{1/8} \xmapsto{} 3^{1/8}\omega^6$
        \item[32. ] $\sigma_{32}: \omega \xmapsto{} \omega^7,\qquad 3^{1/8} \xmapsto{} 3^{1/8}\omega^7$
    
     \end{enumerate}
         \item[\bf{Soln 6}]: We have, $[\mathbb{Q}(2^{1/8},i):\mathbb{Q}(i)]=8$ since $2^{1/8}$ satisfies the polynomial $x^8 -2$ over $\mathbb{Q}(i)$. Also $x^8-2$ is irreducible over $\mathbb{Q}(i)$ as well as separable over $\mathbb{Q}(i)$. So we get that $\mathbb{Q}(2^{1/8},i)$ is the splitting field of $x^8-2$ over $\mathbb{Q}(i)$ , which means that $\mathbb{Q}(2^{1/8},i)/\mathbb{Q}(i)$ is a Galois extension and thus  $|Gal(\mathbb{Q}(2^{1/8},i)/\mathbb{Q}(i))|=[\mathbb{Q}(2^{1/8},i):\mathbb{Q}(i)]=8$. So , from all the options, the correct one is (b)$\mathbb{Z}_8$
         \item[\bf{Soln 7: }] Here, \[
         \begin{aligned}
             [\mathbb{Q}(2^{1/8},i):\mathbb{Q}(\sqrt2)]&=[\mathbb{Q}(2^{1/8},i):\mathbb{Q}(2^{1/8})]\times [\mathbb{Q}(2^{1/8}):\mathbb{Q}(\sqrt2)]\\
             &=2\times4=8
         \end{aligned}
         \] 
         Also, $\mathbb{Q}(2^{1/8},i)$ is the splitting field of $x^4-\sqrt2$ which is separable, irreducible over $\mathbb{Q}(\sqrt2)$, so it is again a Galois extension which means the correct option is (c).
        \item[\bf{soln 8}]: Let $\alpha=\sqrt3+\sqrt7$
    $\begin{aligned}[t]
        &\implies \alpha^2=(\sqrt3+\sqrt7)^2\\
        &\implies \alpha^2=10+2\sqrt21\\ 
        &\implies \alpha^2-10=2\sqrt21\\
        &\implies \alpha^4-20\alpha^2+100=84\\
        &\implies \alpha^4-20\alpha^2+16=0
    \end{aligned}$\\
    So the reqd minimal polynomial is $p(x)=x^4-20x^2+16$
     \item[\bf{Soln 9: }] Option (b) and (d)
    \item[\bf{Soln 10: }] $p(x)=x^4-2x^2-2$ is irreducible by Eisenstein, $p=2$ 
\end{enumerate}

\section{Problem Set 6}

\begin{enumerate}
    \item[\bf{Soln 1: }]  Consider the tower $\mathbb{Q}\subset\mathbb{Q}(2^{1/3})\subset\mathbb{Q}(2^{1/3}, \omega)$. Here $\mathbb{Q}(2^{1/3})$ is not the splitting field of any polynomial, so here we have $B=\mathbb{Q}(2^{1/3})$. Here the minimal polynomial of $2^{1/3}$ is $x^3-2$ whose other roots are $2^{1/3}\omega$ and $2^{1/3}\omega^2$ which are not inside $\mathbb{Q}(2^{1/3})$. So the only automorphism of $\mathbb{Q}(2^{1/3})$ which would send $2^{1/3}$ to another root of its minimal polynomial which is inside $\mathbb{Q}(2^{1/3})$ would be the identity itself. So, $Gal (\mathbb{Q}(2^{1/3})/\mathbb{Q})=\{Id\}$. Now if we take a permutation $\sigma$ from $Gal(\mathbb{Q}(2^{1/3},\omega)/\mathbb{Q})$, then $\sigma$ will fix all elements of the base field $\mathbb{Q}$. So $\sigma\vert_\mathbb{Q}=id$. Now this same $\sigma$ when restricted to $\mathbb{Q}(2^{1/3})$,i.e $\sigma|_{\mathbb{Q}(2^{1/3})}$ is NOT the identity mapping in $Gal(\mathbb{Q}(2^{1/3},\omega)/\mathbb{Q})$. So when we restrict it to $\mathbb{Q}(2^{1/3})$, for it to be in $Gal(\mathbb{Q}(2^{1/3})/\mathbb{Q})$, it should be an automorphism of $\mathbb{Q}(2^{1/3})$ which fixes $\mathbb{Q}$. So $\sigma\vert_{\mathbb{Q}(2^{1/3})}\notin Gal(\mathbb{Q}(2^{1/3})/\mathbb{Q})$. . For example, take the map $\sigma\vert_{\mathbb{Q}(2^{1/3})}: 2^{1/3} \longmapsto2^{1/3}\omega$. This same $\sigma$ is not inside $Gal(\mathbb{Q}(2^{1/3})/\mathbb{Q})$. So, this map $\sigma\vert_{\mathbb{Q}(2^{1/3})}:\mathbb{Q}(2^{1/3}) \longrightarrow \mathbb{Q}(2^{1/3})$ is not a valid map, since the image is not inside $\mathbb{Q}(2^{1/3})$, so it is not an automorphism of $\mathbb{Q}(2^{1/3})$ and thus it is not inside $Gal(\mathbb{Q}(2^{1/3})/\mathbb{Q})$.

    \item[\bf{Soln 2: }] In class, we did the fact that if we drop the condition that $B$ is a splitting field, not that if we drop that $E$ is a splitting field. So I'm not sure if the question is correct. But for the moment, dropping the hypothesis that $B$ is a splitting field, then we can use Soln 1, because the map $\psi(\sigma)=\sigma\vert_B$ would not even be well defined, since $\sigma\vert_B\notin Gal(B/F)$
    \item[\bf{Soln 3: }] option (c). The different maps are :

    \begin{enumerate}
    \item[1. ] $\sigma_1=id$
    \item[2. ] $\sigma_2: \omega \xmapsto{} \omega,\qquad 2^{1/3} \xmapsto{} 2^{1/3}\omega$
    \item[3. ] $\sigma_3: \omega \xmapsto{} \omega,\qquad 2^{1/3} \xmapsto{} 2^{1/3}\omega^2$
    \item[4 ] $\sigma_4: \omega \xmapsto{} \omega^2,\qquad 2^{1/3} \xmapsto{} 2^{1/3}\omega$
    \item[5. ] $\sigma_5: \omega \xmapsto{} \omega^2,\qquad 2^{1/3} \xmapsto{} 2^{1/3}$
    \item[6. ] $\sigma_6: \omega \xmapsto{} \omega^2,\qquad 2^{1/3} \xmapsto{} 2^{1/3}\omega^2$
    \end{enumerate}
    \item[\bf{Soln 4: }] False. Take the tower $\mathbb{Q}\subset\mathbb{Q}(\sqrt2)\subset\mathbb{Q}(2^{1/4})$. Here, $\mathbb{Q}(\sqrt2)/\mathbb{Q}$ is normal and $\mathbb{Q}(2^{1/4})/\mathbb{Q}(\sqrt2)$ is normal but $\mathbb{Q}(2^{1/4})/\mathbb{Q}$ is not normal 
    \item[\bf{Soln 6: }] I think its option (c). The compositum of two fields $E_1$ and $E_2$ is defined to be the intersection of all subfields of the extension $E$ which contain both $E_1$ and $E_2$. So from the given options the compositum of $\mathbb{Q}(2^{1/4},i)$ and $\mathbb{Q}(2^{1/8})$ is the intersection of the of all subfields which contain both of these fields too, so it has to be $\mathbb{Q}(2^{1/8},i)$ 
    \item[\bf{Soln 7:}] The given extension is Galois, so the fixed field $E^G=\mathbb{Q}$. In other words option (c)
    \newpage
    \[
     \begin{aligned}
        &x^4+9=0 \\
        &\implies (x^2)^2+3^2=0 \\
        &\implies (x^2)^2 +\sqrt6x^2+3^2-\sqrt6x^2=0 \\
        &\implies (x^2+3)^2-(\sqrt6x)^2=0\\
        &\implies (x^2-\sqrt6x+3)(x^2+\sqrt6x+3)=0\\
        &\implies x=\dfrac{\sqrt6}{2}(1\pm i), \dfrac{\sqrt6}{2}(-1\pm i)
    \end{aligned}
\]
    Therefore the splitting field is $\mathbb{Q}(\sqrt2,\sqrt3,i)$. The degree is, using tower law,

    \[
    \begin{aligned}
        [\mathbb{Q}(\sqrt2,\sqrt3,i):\mathbb{Q}]=&[\mathbb{Q}(\sqrt3,\sqrt2,i):\mathbb{Q}(\sqrt2,\sqrt3)]\times[\mathbb{Q}(\sqrt2,\sqrt3):\mathbb{Q}(\sqrt2)]\times[\mathbb{Q}(\sqrt2):\mathbb{Q}(i)]\\
        =&2\times2\times2=8
    \end{aligned}
    \]
\end{enumerate}
$\mathcal{DF}$
    
\end{document}